% Source: https://artofproblemsolving.com/community/c3h2389477_dogwater_problems
% Stable identity: aops-topic-2389477-q01
\textbf{Solving the congruences.} From $n\equiv1\pmod7$, $n\in\{1,8,15,\dots\}$; the first of these that is $\equiv3\pmod5$ is $8$, so the first two conditions mean $n\equiv8\pmod{35}$ (by the Chinese remainder theorem, as $5$ and $7$ are coprime). Write $n=8+35t$. Since $35\equiv1\pmod{17}$, the third condition becomes $8+t\equiv11\pmod{17}$, i.e. $t\equiv3\pmod{17}$. So the three conditions together mean
\[n\equiv8+35\cdot3=113\pmod{595},\]
where $595=5\cdot7\cdot17$. (Check: $113=5\cdot22+3=7\cdot16+1=17\cdot6+11$.)

\textbf{Counting.} The lucky positions are $113,\ 113+595,\ 113+2\cdot595,\dots$; the $k$-th is $113+595(k-1)$. The $12$-th is $113+595\cdot11=6658$ and the $13$-th is $113+595\cdot12=7253$. Exactly $12$ people get dogwater when the line has at least $6658$ and at most $7252$ people, so the greatest possible length is $\boxed{7252}$.
